\documentclass[10pt,a4paper]{amsart}
\usepackage[margin=19mm]{geometry}
\usepackage{amsmath,amssymb,mathtools}
\usepackage[scr=rsfs,cal=euler]{mathalfa}
\usepackage{xfrac}
\usepackage[unicode,hidelinks,bookmarks=false]{hyperref}
\newcommand{\ie}{i.e.\ }
\newcommand{\cf}{cf.\ }
\newcommand{\resp}{respectively\ }
\newcommand{\Subsection}{\subsection}
\providecommand{\nobreakdash}{\nobreak\hbox{-}}
\setlength{\emergencystretch}{2em}
\multlinegap0pt
\usepackage{bm}
\usepackage{bbm}
\usepackage{dsfont}
\usepackage{xspace}
\usepackage{cancel}
\usepackage{mathtools}
\usepackage{amsthm}
\makeatletter
\@ifundefined{lem}{\newtheorem{lem}{Lem}}{}
\@ifundefined{prop}{\newtheorem{prop}{Prop}}{}
\makeatother
\makeatletter
\expandafter\ifx\csname centredpic\endcsname\relax
\newenvironment{centredpic}%
{\begin{array}{c}
\setlength{\unitlength}{1em}}%
{\end{array}}
\fi
\expandafter\ifx\csname cproof\endcsname\relax
\newenvironment{cproof}{\noindent\textit{Proof of Claim.}}{\phantom{!}\!\hfill$\diamondsuit$}
\fi
\expandafter\ifx\csname customcor\endcsname\relax
\newenvironment{customcor}[1]
  {\renewcommand\theinnercustomcor{#1}\innercustomcor}
\fi
\expandafter\ifx\csname customthm\endcsname\relax
\newenvironment{customthm}[1]
  {\renewcommand\theinnercustomthm{#1}\innercustomthm}
  {\endinnercustomthm}
\fi
\renewcommand{\ge}{\geqslant}
\renewcommand{\geq}{\geqslant}
\newcommand{\se}{\subseteq}
\newcommand{\sm}{\setminus}
\makeatother
\title[Canonical Decompositions of 3-Connected Graphs]{Canonical Decompositions of 3-Connected Graphs}
\author{Johannes Carmesin, Jan Kurkofka}
\date{Source: arXiv:2304.00945v3 (CC BY 4.0)}
\begin{document}
\maketitle

\noindent\emph{Source and licence.} Canonical Decompositions of 3-Connected Graphs. Version arXiv:2304.00945v3 (CC BY 4.0), distributed under the Creative Commons Attribution 4.0 International licence, as recorded in the arXiv licence metadata for this exact version and shown on the abstract page https://arxiv.org/abs/2304.00945v3. Full rights notice: licenses/arxiv-2304.00945-CC-BY-4.0.txt.\par\smallskip

\noindent\textit{Editorial scope.} Counted unit: the Prop whose source label is \texttt{tetraXtrisep} in the article (source0.tex lines 961--969, source label \texttt{tetraXtrisep}), delivered together with the article's own complete original proof of it (source0.tex lines 971--1010, one \texttt{proof} environment of 40 lines). No mathematics is written, repaired or re-derived: statement and proof are the article's own text line for line. Required context retained uncounted: trivial (lines 855--861); reduceToStrong (lines 947--951). Every definition, display and label of those lines is retained unchanged. Omitted from this package and not redistributed: 1--854, 862--946, 952--960, 970--970, 1011--4234. Nothing inside a retained excerpt is omitted or altered: every retained body is delivered exactly as the article's lines are written, so no omission receipt of any kind accompanies this package. Citations: the retained excerpts carry no citation command at all, so no bibliography entry of the article is transcribed and no reference is redistributed. Reference adaptation: none was needed and none was made; every reference inside the delivered unit resolves inside it, so no unresolved reference or citation is produced. Printed slips kept literal: no printed word, symbol, hypothesis or number of the counted statement or of its proof was altered. Layout and class: the article's own class is \texttt{aic} with packages including epsfig, graphicx, amsbsy, amsmath, amsfonts, amssymb, textcomp, hyperref, aliascnt, mathrsfs, amsthm, etoolbox, tocloft, verbatim; the wrapper is the lane's pinned \texttt{amsart} editorial wrapper at 10pt on A4, which restates the theorem-like environments the retained text needs and adds the article's own macro definitions that the retained text needs (\texttt{\textbackslash ge}, \texttt{\textbackslash geq}, \texttt{\textbackslash se}, \texttt{\textbackslash sm}), with extra wrapper packages (bm, bbm, dsfont, xspace, cancel, mathtools). The delivered numbering is the unit's own. Assets: no retained span contains a figure environment, a graphics-inclusion command, a PDF-page inclusion, a table or a TikZ picture; no image file is copied into this package, the package declares no asset dependency and no image file is redistributed with it. Licence: version arXiv:2304.00945v3 of this article is distributed under the Creative Commons Attribution 4.0 International licence, as recorded in the arXiv licence metadata for this exact version and shown on the abstract page https://arxiv.org/abs/2304.00945v3. A full rights notice is delivered at licenses/arxiv-2304.00945-CC-BY-4.0.txt. Characters: every retained excerpt is pure ASCII, so the delivered engine, XeLaTeX with its default Latin Modern text font, reports no missing character.
\begin{lem}\label{trivial}
The following are equivalent for every tri-separation $(A,B)$ of a 3-connected graph~$G$:
 \begin{enumerate}
  \item $(A,B)$ is trivial;
  \item $A$ and $B$ are the two sides of an atomic cut.
 \end{enumerate}
\end{lem}

\begin{lem}\label{reduceToStrong}
Let $(A,B)$ be a mixed 3-separation of a 3-connected graph~$G$.
For every edge $uv$ in~$G$ with both ends $u,v$ in the separator of $(A,B)$,
there is a strengthening $(A',B')$ of $(A,B)$ so that $u,v\in B'$.
\end{lem}

\begin{prop}\label{tetraXtrisep}
For every 3-connected graph~$G$, the following assertions are equivalent:
\begin{enumerate}
    \item\label{ItemTetra} $G$ is internally 4-connected or~$G\in\{K_4,K_{3,3}\}$;
        \item\label{ItemTetra_neu} every 3-separation of $G$ is trivial or~$G=K_4$;
        \item\label{ItemTrivial1} all tri-separations of~$G$ are trivial or $G=K_4$;
    \item\label{ItemTrivial2} all strong tri-separations of~$G$ are trivial.
\end{enumerate}
\end{prop}

\begin{proof}
\ref{ItemTetra}$\rightarrow$\ref{ItemTetra_neu}.
If $G$ is internally 4-connected or $G=K_{3,3}$, then every 3-separation $(A,B)$ of $G$ has a side ($A$, say) that induces a claw; in particular, $G[A]$ contains no cycle, so $(A,B)$ is trivial.

$\neg$\ref{ItemTrivial1}$\to\neg$\ref{ItemTetra_neu}.
Let $(A,B)$ be a nontrivial tri-separation of $G$.
Recall that $A\sm B$ and $B\sm A$ are nonempty by the definition of mixed-separation.
For each edge in $S(A,B)$ we pick one of its endvertices and add it to both sides. 
We pick these endvertices so that we preserve that $A\sm B$ and $B\sm A$ are nonempty. 
This is possible as $G$ has at least five vertices. As this preserves nontriviality, we end up with a nontrivial 3-separation of $G$. 

Clearly \ref{ItemTrivial1}$\to $\ref{ItemTrivial2}.



$\neg$\ref{ItemTetra}$\to\neg$\ref{ItemTrivial2}.
As $G$ is not internally 4-connected and $G\notin\{K_4,K_{3,3}\}$, we may let $(A,B)$ be a 3-separation of~$G$ none of whose sides induces a claw.
Let~$X:=A\cap B$.
If $G$ had at most four vertices, then~$G$ would be a~$K_4$, contradicting our assumption, so we have~$|G|\ge 5$.

{\bf Case 1:} the induced subgraph $G[X]$ has no edges.
Then $|A\sm B|\geq 2$ and~$|B\sm A|\geq 2$.
By \autoref{reduceToStrong}, there is a strong tri-separation $(A',B')$ of~$(A,B)$ with $A\sm B\se A'$ and $B\sm A\se B'$.
Since $A\sm B$ and $B\sm A$ have size at least two, so have $A'$ and~$B'$.
Thus~$(A',B')$ is nontrivial by \autoref{trivial}.

{\bf Case 2:} the induced subgraph $G[X]$ contains an edge~$x_1 x_2$.
Recall that $|G|\ge 5$.
The only 3-connected graphs on exactly five vertices that have a 3-separator are the 4-wheel and
$K_5^-$ ($K_5$~minus one edge). The unique $3$-separator of
$K_5^-$ is a triangle and thus it is the separator of a nontrivial tri-separation.
The 4-wheel has two nontrivial strong tri-separations.
So it remains to consider the
case that $G$ has at least six vertices. 
By symmetry, we may assume that $|A\sm B|\geq 2$.
Since $x_1 x_2$ is an edge of~$G$ with both ends in $X=A\cap B$, we find a strong tri-separation $(A',B')$ of~$G$ with $A\sm B\se A'$ and $B\sm A\se B'$ such that $x_1,x_2\in B'$.
Hence, the side~$B'$ misses at most one vertex of~$B$.
As $|B|\geq 4$, this gives~$|B'|\geq 3$. 
We also have~$|A'|\geq |A\sm B|$, and $|A\sm B|\ge 2$ by assumption. Hence the strong tri-separation $(A',B')$ is nontrivial by \autoref{trivial}.
\end{proof}

\end{document}
